\section{Same-Target Temporal Probe}
\label{sec:same-target}

At control step $t$, a Fresh component uses offset 0 from the chunk predicted
from the current observation. A component with source age $d$ instead uses
offset $d$ from the chunk queried $d$ steps earlier. Both entries were
originally predicted for the same physical step $t$. For example, A0G20 takes
the arm command from query $t$ at offset 0 and the gripper command from query
$t-20$ at offset 20. A20G0 swaps these assignments: the arm comes from query
$t-20$ at offset 20, and the gripper comes from query $t$ at offset 0. During
the first 20 steps, before the older source exists, every component executes
the Fresh offset-0 command.

Formally, let $q$ be policy-query time, $k$ the chunk offset, and $A_q[k]$
the predicted action entry. A valid same-target source satisfies
\begin{equation}
q + k = t .
\label{eq:same-target}
\end{equation}
The source age is $d=t-q$, so same-target alignment requires $k=d$. We
therefore measure the joint effect of an older source observation and a
longer-lookahead prediction; the probe does not separate those factors.
Translation comprises action dimensions 0--2, rotation dimensions 3--5, and
the gripper dimension 6. Same-target probes query the policy once per executed
step, giving a policy-query rate of 1.0, and discard unselected predictions.
They are diagnostic interventions, not deployment executors.

Executable schedules use separate notation. Periodic H$h$ queries every $h$
steps and commits the full action to that chunk, giving the frozen H16, H4, and
H2 schedules. ARM$h$\_GRIP$g$ assigns separate commitment horizons to the arm
and gripper, yielding ARM4\_GRIP32 and ARM2\_GRIP16. TE\_DENSE queries at every
step and applies the canonical temporal ensemble. These frozen schedules test
an operational consequence; they are not a new adaptive executor.
